\section{Aplicación móvil}
\label{sec:movil}

\subsection{Justificación cuantitativa de Android nativo}

El Módulo C permite dos rutas para el cliente móvil obligatorio: Android nativo (Kotlin +
Jetpack Compose) o multiplataforma (Flutter 3.x o React Native 0.74+). El Anexo B de la guía de
esta entrega exige justificar esa elección con al menos dos criterios cuantitativos, no solo
preferencia estética —conforme a lo discutido en la literatura sobre desarrollo de aplicaciones
móviles~\cite{wasserman2010mobile,elkassas2017taxonomy,biornhansen2019crossplatform}. El cliente
móvil de este proyecto es para técnicos de campo: usa cámara y geolocalización de forma continua
durante la jornada, en dispositivos de gama variable asignados por la operadora. Se eligió
\textbf{Android nativo con Kotlin + Jetpack Compose} (MVVM, inyección de dependencias manual),
sustentado en dos mediciones directas sobre el proyecto ya construido (ADR-0006):

\begin{enumerate}[nosep]
  \item \textbf{Tamaño del artefacto instalable}: el APK de depuración pesa \textbf{11.06\,MB},
        sin ningún motor de renderizado embebido adicional —Jetpack Compose compila a las mismas
        vistas nativas de Android (\emph{bytecode} ART), a diferencia de un framework
        multiplataforma que empaqueta su propio motor (Skia/Flutter engine, o el puente
        JavaScript de React Native) dentro de cada instalación.
  \item \textbf{Dependencias de terceros para las capacidades obligatorias}: cámara y
        geolocalización se resuelven con \textbf{cero dependencias de interoperabilidad} —la
        cámara usa \texttt{ActivityResultContracts.TakePicture} (parte del propio SDK de Android)
        y la geolocalización una única dependencia oficial de Google
        (\texttt{play-services-location}). Un framework multiplataforma necesitaría, en cambio,
        un \emph{plugin} puente con la sobrecarga de serialización y el riesgo de mantenimiento
        que eso implica.
\end{enumerate}

Como criterio cualitativo adicional: el backend del proyecto ya es JVM (Java 21 + Spring Boot),
por lo que Kotlin comparte máquina virtual y herramientas de \emph{build} (Gradle) con el resto
del equipo, reduciendo el salto conceptual frente a adoptar Dart o un puente nativo de JavaScript
desde cero.

\subsection{Arquitectura MVVM y modo sin conexión}

Cada pantalla sigue el patrón \textbf{MVVM}: un \texttt{ViewModel} concentra el estado
(\texttt{StateFlow}) y la lógica, mientras la vista en Jetpack Compose solo observa y dibuja. El
listado de tickets asignados usa una caché local con Room, de forma que un técnico en campo sin
cobertura de red siga viendo su carga de trabajo del día.

\subsection{Capacidades del dispositivo}

Se integran las dos capacidades obligatorias del Módulo C:

\begin{itemize}[nosep]
  \item \textbf{Cámara}: evidencia fotográfica del cierre de un ticket, capturada con
        \texttt{ActivityResultContracts.TakePicture} y un \texttt{FileProvider} propio de la
        aplicación (sin escribir a almacenamiento externo compartido).
  \item \textbf{Geolocalización}: ubicación del cierre capturada con
        \texttt{FusedLocationProviderClient}, solicitada solo en el momento del cierre, no como
        rastreo continuo en segundo plano, y precedida de un diálogo explícito de consentimiento
        que declara la finalidad y el destino del dato antes de pedir el permiso del sistema
        operativo (\texttt{TicketDetailScreen.kt}).
\end{itemize}

La foto y las coordenadas no se quedan en el teléfono: \texttt{TicketRepository.closeOnSite}
las envía junto con el cambio de estado a \texttt{RESUELTO} en una sola llamada
\texttt{PATCH /api/v1/tickets/\{id\}/status} (foto codificada en Base64), con reintento —hasta 3
veces ante \texttt{IOException}, con espera de 1 y 2\,s— para tolerar una red de campo
inestable. El servicio principal las guarda como columnas \texttt{BYTES}/\texttt{FLOAT8}
adicionales de la fila del ticket (migración \texttt{V3\_\_add\_close\_evidence.sql}), sin un
almacén de objetos aparte. El reintento es idempotente: si el primer \texttt{PATCH} sí llegó al
servidor y solo se perdió la respuesta, el segundo no vuelve a mover \texttt{resolvedAt} ni a
recalcular el incumplimiento de SLA, y el servidor exige la evidencia cuando quien cierra es un
\texttt{TECNICO} (un \texttt{ADMIN} puede seguir resolviendo un ticket sin pasar por el móvil).

La Figura~\ref{fig:movil-capturas} muestra ambas capacidades reales en el emulador
(\texttt{Medium\_Phone\_API\_36.0}), instalando el mismo APK de depuración que produce el
\emph{job} \texttt{build-mobile-apk} de CI/CD: el listado (que además exhibe en vivo el modo sin
conexión de la Sección anterior --el emulador arrancó con datos ya cacheados en Room, y la
aplicación lo declara explícitamente en vez de mostrar una lista vacía o congelarse) y la pantalla
de detalle con los botones reales \texttt{Tomar Foto} y \texttt{GPS Cierre}.

\begin{figure}[H]
\centering
\begin{minipage}{0.47\textwidth}
  \centering
  \includegraphics[width=\linewidth]{../../release/screenshots/mobile_01_lista.png}
  \\ \footnotesize (a) Listado (modo sin conexión, datos cacheados con Room)
\end{minipage}
\hfill
\begin{minipage}{0.47\textwidth}
  \centering
  \includegraphics[width=\linewidth]{../../release/screenshots/mobile_02_detalle_orden.png}
  \\ \footnotesize (b) Detalle: cámara y GPS de cierre
\end{minipage}
\caption{Capturas reales de \texttt{apps/mobile} en emulador Android (2026-09-03).}
\label{fig:movil-capturas}
\end{figure}

La pantalla de detalle muestra además el identificador completo del ticket en un elemento de
texto seleccionable —una decisión deliberada para poder demostrar, en una presentación en vivo,
que la aplicación móvil y la aplicación web leen y escriben sobre el mismo backend real: cerrar
un ticket desde el teléfono y encontrarlo actualizado en la consola web buscando por ese mismo
identificador (Sección~\ref{sec:web}).

\subsection{Pruebas}

La capa de \texttt{ViewModel} tiene pruebas unitarias con JUnit 5 y \texttt{kotlinx-coroutines-test}
(\texttt{LoginViewModelTest}). La prueba instrumentada de extremo a extremo exigida por el Módulo
C (Espresso o Compose Testing) cubre el flujo de validación de formulario de inicio de sesión
(\texttt{LoginScreenTest}, Compose Testing) contra la actividad real de la aplicación, sin
depender de que el backend esté disponible —de forma estable en un pipeline de integración
continua.
